
\documentclass{article}
\usepackage{colortbl}
\usepackage{makecell}
\usepackage{multirow}
\usepackage{supertabular}

\begin{document}

\newcounter{utterance}

\centering \large Interaction Transcript for game `bbh', experiment `disambiguation\_qa', episode 180 with reflexion{-}Qwen3.5{-}9B.
\vspace{24pt}

{ \footnotesize  \setcounter{utterance}{1}
\setlength{\tabcolsep}{0pt}
\begin{supertabular}{c@{$\;$}|p{.15\linewidth}@{}p{.15\linewidth}p{.15\linewidth}p{.15\linewidth}p{.15\linewidth}p{.15\linewidth}}
    \# & \multicolumn{2}{c}{Player} && \multicolumn{2}{c}{Game Master} \\
    \hline

    \theutterance \stepcounter{utterance}  
    & & & \multicolumn{4}{p{0.6\linewidth}}{
        \cellcolor[rgb]{0.9,0.9,0.9}{
            \makecell[{{p{\linewidth}}}]{
                \texttt{\tiny{[P1$\langle$GM]}}
                \texttt{Clarify the meaning of sentences with ambiguous pronouns.} \\
\\ 
\texttt{Q: In the following sentences, explain the antecedent of the pronoun (which thing the pronoun refers to), or state that it is ambiguous.} \\
\texttt{Sentence: The chief told the counselor that they took the day off.} \\
\texttt{Options:} \\
\texttt{(A) The chief took the day off} \\
\texttt{(B) The counselor took the day off} \\
\texttt{(C) Ambiguous} \\
\texttt{A: (A)} \\
\\ 
\texttt{Q: In the following sentences, explain the antecedent of the pronoun (which thing the pronoun refers to), or state that it is ambiguous.} \\
\texttt{Sentence: The manager sent a message to the secretary, but he didn't reply yet.} \\
\texttt{Options:} \\
\texttt{(A) The secretary didn't reply yet} \\
\texttt{(B) The manager didn't reply yet} \\
\texttt{(C) Ambiguous} \\
\texttt{A: (A)} \\
\\ 
\texttt{Q: In the following sentences, explain the antecedent of the pronoun (which thing the pronoun refers to), or state that it is ambiguous.} \\
\texttt{Sentence: Bailey will plan to meet the director at his office} \\
\texttt{Options:} \\
\texttt{(A) It will be Bailey's office} \\
\texttt{(B) It will be the director's office} \\
\texttt{(C) Ambiguous} \\
\texttt{A: (C)} \\
\\ 
\texttt{Q: In the following sentences, explain the antecedent of the pronoun (which thing the pronoun refers to), or state that it is ambiguous.} \\
\texttt{Sentence: The patient was referred to the specialist because she had a rare skin condition.} \\
\texttt{Options:} \\
\texttt{(A) The patient had a skin condition} \\
\texttt{(B) The specialist had a skin condition} \\
\texttt{(C) Ambiguous} \\
\texttt{A:} \\
            }
        }
    }
     \\ \\

    \theutterance \stepcounter{utterance}  
    & \multicolumn{4}{p{0.6\linewidth}}{
        \cellcolor[rgb]{0.561,0.737,0.561}{
            \makecell[{{p{\linewidth}}}]{
                \texttt{\tiny{[P1$\rangle$GM]}}
                \texttt{(A)} \\
\\ 
\texttt{**Reasoning:**} \\
\texttt{In this sentence, the pronoun "she" refers to the person who possesses the condition that necessitated the referral. Since the patient is the one being referred *because* of the condition, it logically follows that the patient is the one with the rare skin condition, not the specialist. Therefore, the antecedent is clearly "The patient."} \\
            }
        }
    }
    & & \\ \\

    \theutterance \stepcounter{utterance}  
    & & & \multicolumn{2}{p{0.3\linewidth}}{
        \cellcolor[rgb]{0.9,0.9,0.9}{
            \makecell[{{p{\linewidth}}}]{
                \texttt{\tiny{[GM$|$GM]}}
                \texttt{Target: (A)} \\
            }
        }
    }
    & & \\ \\

    \theutterance \stepcounter{utterance}  
    & & & \multicolumn{2}{p{0.3\linewidth}}{
        \cellcolor[rgb]{0.9,0.9,0.9}{
            \makecell[{{p{\linewidth}}}]{
                \texttt{\tiny{[GM$|$GM]}}
                \texttt{game\_result = LOSE} \\
            }
        }
    }
    & & \\ \\

    \theutterance \stepcounter{utterance}  
    & & \multicolumn{4}{p{0.6\linewidth}}{
        \cellcolor[rgb]{0.95,0.95,0.95}{
            \makecell[{{p{\linewidth}}}]{
                \texttt{\tiny{[P1$\leftrightarrow$P]}}
                \texttt{Post{-}episode Reflexion (not sent to GM):} \\
\texttt{New plan: I incorrectly assumed that the pronoun "she" must refer to the patient because the patient is the subject of the first clause, ignoring the possibility that the specialist could also be female and possess the condition. In ambiguous pronoun resolution tasks, I should not rely solely on grammatical proximity or narrative logic; instead, I must explicitly evaluate if the sentence structure allows for multiple valid antecedents. Next time, I will analyze whether the verb phrase "had a rare skin condition" logically applies to both potential referents (the patient and the specialist). If both are plausible, I will select option (C) Ambiguous rather than forcing a single interpretation based on assumption.} \\
            }
        }
    }
    & \\ \\

\end{supertabular}
}

\end{document}
